We studied optimizer choice as a geometry intervention for feature-based
uncertainty. In an accuracy-controlled selected-configuration study with SGD,
Adam, and AdamW, we find that comparable ID performance can hide substantial
differences in penultimate representation geometry and detector-family behavior.
Selected optimizers induce different fingerprints in neural-collapse geometry, feature
norms, within-class dispersion, and class separation, and detector families read
these fingerprints differently.

The clearest example is the AdamW-selected raw-distance split. These
configurations can remain competitive under logit-based readouts while showing
large near-OOD gaps under raw Mahalanobis and kNN. L2 normalization recovers much
of this gap, implicating a radial or norm-coupled failure channel, whereas
covariance shrinkage has a more limited and detector-dependent effect. These
results support a compatibility view rather than an optimizer-good/bad ranking.

The key lesson is that optimizer, penultimate representation geometry, and
detector readout should be evaluated jointly. Feature-based uncertainty should
therefore be evaluated as a geometry--detector compatibility problem, not as an
accuracy-only model-selection problem.
